\documentclass[11pt]{article}

% ============================================================================
%  Master Chess — Field Manual
%  Typeset with Tectonic (XeTeX).  Build:  tectonic master-chess-manual.tex
% ============================================================================

\usepackage{fontspec}
\usepackage{amsmath}
\usepackage[a4paper,margin=2.4cm,headsep=18pt,footskip=26pt]{geometry}
\usepackage{xcolor}
\usepackage{fancyhdr}
\usepackage{titlesec}
\usepackage{tcolorbox}
\usepackage{enumitem}
\usepackage{booktabs}
\usepackage{array}
\usepackage{colortbl}
\usepackage{tabularx}
\usepackage{xltabular}
\usepackage{graphicx}
\usepackage{pifont}
\usepackage[hidelinks]{hyperref}
\usepackage{needspace}
\tcbuselibrary{skins,breakable}
\usetikzlibrary{calc}

% ---- Palette: a wooden chess set + score classifications ------------------
\definecolor{walnut}{HTML}{3A2A1E}     % dark board wood — headings, rules
\definecolor{walnutmid}{HTML}{5C4632}  % mid wood
\definecolor{parchment}{HTML}{F4ECD8}  % board-light square / page tint
\definecolor{parchdeep}{HTML}{E8DCC0}  % deeper parchment for rules
\definecolor{forest}{HTML}{2E5D45}     % primary accent (felt green)
\definecolor{forestlt}{HTML}{4A7C4E}   % "best" green
\definecolor{blunder}{HTML}{B23A2E}    % blunder red
\definecolor{gold}{HTML}{C9A227}       % highlight gold
\definecolor{ink}{HTML}{241B14}        % body text
\definecolor{slate}{HTML}{6B5D4F}      % muted captions
\definecolor{rulehair}{HTML}{C9B89A}   % hairline

% ---- Fonts ----------------------------------------------------------------
\setmainfont{Charter}[
  Numbers = OldStyle,
  Ligatures = TeX ]
\setsansfont{DejaVu Sans}[
  Scale = 0.92,
  LetterSpace = 0.4 ]
\setmonofont{DejaVu Sans Mono}[Scale = 0.82]
\newfontfamily\tabnum{Charter}[Numbers = Lining]
\newfontfamily\chessfont{DejaVu Sans}
\newcommand{\king}{{\chessfont\symbol{"2654}}}   % ♔
\newcommand{\yes}{{\color{forestlt}\ding{52}}}    % ✓ supports
\newcommand{\no}{{\color{blunder}\ding{56}}}      % ✗ counts against

\color{ink}
\linespread{1.10}
\setlength{\parindent}{0pt}
\setlength{\parskip}{6pt}

% ---- Headings -------------------------------------------------------------
\titleformat{\section}
  {\needspace{4\baselineskip}\sffamily\bfseries\Large\color{walnut}}
  {\colorbox{walnut}{\textcolor{parchment}{\;\thesection\;}}}{0.7em}{}
  [\vspace{-6pt}{\color{gold}\rule{\linewidth}{1.4pt}}]
\titleformat{\subsection}
  {\needspace{3\baselineskip}\sffamily\bfseries\large\color{forest}}
  {\thesubsection}{0.6em}{}
\titleformat{\subsubsection}
  {\needspace{2\baselineskip}\sffamily\bfseries\normalsize\color{walnutmid}}
  {}{0em}{}
\titlespacing*{\section}{0pt}{18pt}{8pt}
\titlespacing*{\subsection}{0pt}{12pt}{4pt}
\titlespacing*{\subsubsection}{0pt}{8pt}{2pt}

% ---- Header / footer ------------------------------------------------------
\pagestyle{fancy}
\fancyhf{}
\renewcommand{\headrulewidth}{0pt}
\renewcommand{\footrulewidth}{0pt}
\fancyhead[L]{\sffamily\footnotesize\color{slate}\MakeUppercase{Master Chess}}
\fancyhead[R]{\sffamily\footnotesize\color{slate}\leftmark}
\fancyfoot[L]{\sffamily\footnotesize\color{slate}Field Manual}
\fancyfoot[R]{\sffamily\footnotesize\color{slate}\thepage}
\fancyfoot[C]{\footnotesize\color{gold}\ding{169}}
\renewcommand{\sectionmark}[1]{\markboth{\thesection\quad #1}{}}

% ---- Table helpers --------------------------------------------------------
\newcolumntype{Y}{>{\raggedright\arraybackslash}X}
\renewcommand{\arraystretch}{1.28}
\arrayrulecolor{rulehair}

% ---- Boxes ----------------------------------------------------------------
% Example-of-play box (tabletop-manual tradition)
\newtcolorbox{playbox}[1][Example of Play]{
  breakable, enhanced,
  colback=parchment, colframe=walnut,
  boxrule=0pt, leftrule=4pt, arc=2pt, outer arc=2pt,
  left=14pt, right=14pt, top=10pt, bottom=12pt,
  boxsep=2pt,
  overlay unbroken and first={
    \node[anchor=north west, font=\sffamily\bfseries\footnotesize,
          text=parchment, fill=walnut, inner xsep=8pt, inner ysep=3pt]
      at ([xshift=10pt,yshift=10pt]frame.north west) {\MakeUppercase{#1}};},
  before skip=12pt, after skip=12pt }

% Rule / definition box
\newtcolorbox{rulebox}[1]{
  breakable, enhanced,
  colback=parchdeep, colframe=forest,
  boxrule=0pt, leftrule=4pt, arc=1pt, outer arc=1pt,
  left=12pt, right=12pt, top=8pt, bottom=10pt,
  fonttitle=\sffamily\bfseries\small,
  coltitle=parchment, title={\ding{110}\enspace #1},
  before skip=10pt, after skip=10pt }

% Note / caution box
\newtcolorbox{notebox}[1][Note]{
  breakable, enhanced,
  colback=white, colframe=gold,
  boxrule=0pt, leftrule=4pt, arc=1pt,
  left=12pt, right=12pt, top=7pt, bottom=8pt,
  fonttitle=\sffamily\bfseries\footnotesize,
  coltitle=walnutmid, colbacktitle=white,
  attach title to upper={\par\vskip4pt},
  title={\ding{72}\enspace\MakeUppercase{#1}},
  before skip=10pt, after skip=10pt }

% Classification chip
\newcommand{\chip}[2]{\,\colorbox{#1}{\textcolor{white}{\sffamily\bfseries\footnotesize\,#2\,}}\,}
\newcommand{\best}{\chip{forestlt}{Best}}
\newcommand{\good}{\chip{forest}{Good}}
\newcommand{\inacc}{\chip{gold}{Inaccuracy}}
\newcommand{\mistake}{\chip{walnutmid}{Mistake}}
\newcommand{\blun}{\chip{blunder}{Blunder}}

\newcommand{\fen}[1]{{\ttfamily\small #1}}
\newcommand{\code}[1]{{\ttfamily #1}}
\newcommand{\skl}[1]{{\sffamily\small\color{forest}#1}}

\hypersetup{pdftitle={Master Chess — Field Manual},
            pdfauthor={Master Chess},
            pdfsubject={Adaptive educational chess tutor — user \& mechanics manual}}

% ============================================================================
\begin{document}

% ---- Cover (full-bleed via current-page overlay) --------------------------
\begin{titlepage}
\thispagestyle{empty}
\begin{tikzpicture}[remember picture, overlay]
  % Full-bleed walnut background
  \fill[walnut]
    (current page.south west) rectangle (current page.north east);
  % Gold inner frame
  \draw[gold, line width=1.2pt]
    ($(current page.south west)+(1.7,1.7)$) rectangle
    ($(current page.north east)+(-1.7,-1.7)$);
  % Large faint king watermark
  \node[opacity=0.09, text=parchment, scale=28]
    at ($(current page.center)+(0,-1.0)$) {\king};
  % Eyebrow
  \node[anchor=west, text=gold, font=\sffamily\bfseries, scale=1.05]
    at ($(current page.west)+(2.6,5.4)$) {\MakeUppercase{Field Manual}};
  % Title
  \node[anchor=west, text=parchment, font=\sffamily\bfseries, scale=3.6]
    at ($(current page.west)+(2.5,3.7)$) {Master Chess};
  % Rule
  \draw[gold, line width=1pt]
    ($(current page.west)+(2.6,2.5)$) -- ($(current page.west)+(10.2,2.5)$);
  % Tagline
  \node[anchor=west, text=parchment, font=\itshape, scale=1.4]
    at ($(current page.west)+(2.6,1.6)$)
    {Turn your own games into your training plan.};
  % Diagnose / Explain / Prescribe
  \node[anchor=west, text=gold, font=\sffamily\bfseries, scale=1.15]
    at ($(current page.west)+(2.6,-2.0)$)
    {DIAGNOSE \quad\textcolor{parchment}{\ding{117}}\quad
      EXPLAIN \quad\textcolor{parchment}{\ding{117}}\quad PRESCRIBE};
  % Footer block
  \node[anchor=west, text=parchment, font=\sffamily, scale=1.0]
    at ($(current page.south west)+(2.6,4.2)$)
    {An adaptive educational chess tutor};
  \node[anchor=west, text=slate, font=\sffamily, scale=0.9]
    at ($(current page.south west)+(2.6,3.4)$)
    {Player \& Mechanics Manual \textbullet\ Version 1.0};
\end{tikzpicture}
\end{titlepage}

% ---- How to read this book ------------------------------------------------
\thispagestyle{empty}
\vspace*{1.5cm}
{\sffamily\bfseries\LARGE\color{walnut} How to read this manual}\\[2pt]
{\color{gold}\rule{\linewidth}{1.4pt}}

\vspace{6pt}
This manual is written the way a good board-game rulebook is written: the rules
first, then a fully worked \emph{game of play} that walks one real game from
upload to a finished training plan, so you can see every mechanism act on real
moves. Wherever you see a tinted box, it is doing one of three jobs.

\begin{playbox}[Example of Play — how these boxes work]
Boxes like this one follow a single running example — the analysis of a real
historical game — so that every abstract rule is immediately grounded in
something concrete. Read them in order and you will have watched the whole
system operate end to end.
\end{playbox}

\begin{rulebox}{Rule box — the exact mechanic}
Boxes like this state a precise rule: a threshold, a formula, a piece of
scheduling logic. These are the authoritative numbers, taken directly from the
system's source of record.
\end{rulebox}

\begin{notebox}[Note]
Boxes like this flag something easy to get wrong, a limitation, or a
cross-reference worth keeping in mind.
\end{notebox}

\vspace{4pt}
\begin{center}
\begin{tcolorbox}[enhanced, width=0.92\linewidth, colback=parchment,
  colframe=rulehair, boxrule=0.6pt, arc=2pt, left=14pt,right=14pt,top=8pt,bottom=8pt]
{\sffamily\footnotesize\color{walnutmid}\MakeUppercase{The one idea to keep}}\\[2pt]
Master Chess is not an engine with a chat window bolted on. It runs a
\textbf{closed learning loop} — \textsc{diagnose} what you can and cannot do,
\textsc{explain} it against engine truth, \textsc{prescribe} exactly what to
study and drill — and \textbf{every score it shows you has a receipt}: the exact
moves that produced it. No black box.
\end{tcolorbox}
\end{center}

\newpage
\tableofcontents
\newpage

% ===========================================================================
\section{What Master Chess Is}
% ===========================================================================

Master Chess is an adaptive educational chess tutor. You upload your own games;
it analyses every move with the Stockfish engine, and from that
engine-grounded analysis it builds a persistent, \emph{explainable} model of
your ability across a taxonomy of \textbf{27 skills}. It then diagnoses your
playing \textbf{level} and your current \textbf{plateau}, writes you a
personalised \textbf{training plan} with matched reading, and turns the very
mistakes you made into spaced-repetition \textbf{drills}.

The product's thesis is that expert chess knowledge should be encoded into a
structured substrate, and instruction rendered dynamically for the individual
learner from engine truth --- never from an engine's chat-layer guesswork.

\subsection{The two learning loops (and why the AI is neither)}

The system has exactly two places where ``learning'' happens, and a language
model is not one of them:

\begin{enumerate}[leftmargin=1.4em, itemsep=2pt]
  \item \textbf{The Player Model} --- a structured representation of you across
  27 skills, updated by Bayesian Knowledge Tracing from engine-verified
  evidence in every game and drill.
  \item \textbf{The Chess Knowledge Layer} --- engine ground truth, an
  opening / master-game library, and a book corpus mapped to the skill
  taxonomy.
\end{enumerate}

An \textbf{AI reasoning layer} sits on top to write prose explanations, but it
is always subordinate to engine truth: a runtime \emph{claim guard} rejects any
generated sentence that contradicts a verified move or evaluation. If the model
cannot say something true, the app falls back to a deterministic, engine-only
readout rather than let an unverified sentence reach you.

\begin{rulebox}{What grounds every claim}
The engine (Stockfish) and the classified move record are the source of truth.
The player model is derived from them by fixed, inspectable rules. The language
model may only \emph{phrase} facts it is handed; it may never introduce a move,
number, or verdict of its own. Any sentence that fails verification is
discarded.
\end{rulebox}

\subsection{Where the AI is used --- and where it isn't}\label{sec:ai}

Master Chess uses AI deliberately and narrowly. Everything that \emph{scores,
judges, or decides} is deterministic --- the engine and fixed rules. The AI is a
\textbf{language layer} that turns those verified facts into readable coaching,
and answers your questions. It appears in exactly four places:

\begin{center}
\begin{tabularx}{\linewidth}{@{}l Y l@{}}
\toprule
\textbf{Where} & \textbf{What the AI does} & \textbf{Grounded by} \\
\midrule
Game Review summary & Writes the post-game coaching note from your result, accuracy, and worst move & Claim guard + deterministic fallback \\
Prophylaxis notes & Phrases \emph{why} a defensive move mattered & Deterministic probe decides the score; AI only narrates \\
Progress page & Writes your ``coach's summary'' from your real stats & Real numbers only + deterministic fallback \\
Help assistant & Answers ``how do I\dots'' questions in chat & Fixed knowledge base; refuses out-of-scope \\
\bottomrule
\end{tabularx}
\end{center}

In every case the same rule holds: \textbf{the AI may only rephrase what the
engine and your data already establish.} It never sets a mastery number, never
classifies a move, never chooses your next action (a fixed rule does that),
never invents chess advice, never generates a drill hint (those are derived from
the engine's best move), and \textbf{never plays your opponent in Play vs AI ---
that is Stockfish}. If the AI backend is unavailable, every one of these
surfaces degrades gracefully to a deterministic, engine-only version --- the app
keeps working, just in plainer prose. The AI provider is pluggable
(a self-hosted model by default), so no third-party service is required for the
core loop.

\subsection{How it runs}

\begin{center}
\begin{tabularx}{\linewidth}{@{}l Y@{}}
\toprule
\textbf{Layer} & \textbf{Technology} \\
\midrule
Client & React single-page app served by Vite (development port \code{5173}) \\
Server & Express API on \code{tsx} (default port \code{8030}) \\
Storage & SQLite via Drizzle ORM (\code{better-sqlite3}, WAL mode) \\
Engine & Stockfish over UCI, pooled; default search depth \code{16}, multi-PV \code{4} \\
Reasoning & Provider-abstracted LLM (local Ollama by default; hosted models by tier) \\
\bottomrule
\end{tabularx}
\end{center}

To run it locally: \code{npm run dev} starts client and server together; the
server first applies database migrations, seeds the master-game library,
resumes any interrupted analyses, and purges expired sessions. A health probe
lives at \code{GET /api/health}.

\begin{notebox}[Scope note]
This manual documents the \textbf{Master Chess} standalone app (the
\code{ai\_chess} project). A sibling project in the same workspace is a separate
classroom product with teacher/student dashboards; those screens are \emph{not}
part of Master Chess and are not described here.
\end{notebox}

% ===========================================================================
\section{The Vocabulary: Skills, Levels, Plateaus}
% ===========================================================================

Everything the app says about you is expressed in three linked vocabularies.
Learn these and the rest of the manual reads easily.

\subsection{The 27 skills, in 4 categories}

Each skill carries a mastery score from 0--100, a trend arrow, and --- once
there is evidence --- a list of the exact moves that produced the score. Three
skills marked \textit{partial} are only partially diagnosable from game data
alone.

\begin{center}
\tabnum
\begin{xltabular}{\linewidth}{@{}p{1.4em} Y l@{}}
\toprule
\# & \textbf{Skill} & \textbf{Diagnosable} \\
\midrule
\endhead
\multicolumn{3}{@{}l}{\textbf{\color{forest}OPENING}}\\
1 & Opening principles & yes \\
2 & Opening repertoire & yes \\
3 & Opening preparation \& novelties & partial \\
\addlinespace
\multicolumn{3}{@{}l}{\textbf{\color{forest}MIDDLEGAME}}\\
4 & Strategic planning & yes \\
5 & Positional play & yes \\
6 & Patience \& long-term thinking & yes \\
7 & Prophylaxis & yes \\
8 & Tactical pattern recognition & yes \\
9 & Tactical consistency (blunder prevention) & yes \\
10 & Calculation precision & yes \\
11 & Attack \& initiative & yes \\
12 & Piece activity \& coordination & yes \\
13 & Defence \& counterplay & yes \\
14 & Converting advantages & yes \\
\addlinespace
\multicolumn{3}{@{}l}{\textbf{\color{forest}ENDGAME}}\\
15 & Endgame principles & yes \\
16 & Game transition (middlegame to endgame) & yes \\
17 & Endgame precision \& conversion & yes \\
18 & Pawn endings & yes \\
19 & Rook endings & yes \\
20 & Knight endings & yes \\
21 & Bishop \& mixed-piece endings & yes \\
\addlinespace
\multicolumn{3}{@{}l}{\textbf{\color{forest}PSYCHOLOGY \& MENTAL}}\\
22 & Thought process \& candidate moves & partial \\
23 & Pattern recognition speed & yes \\
24 & Intuition \& practical judgment & partial \\
25 & Time management & yes \\
26 & Psychological resilience & yes \\
27 & Psychological warfare & partial \\
\bottomrule
\end{xltabular}
\end{center}

\begin{notebox}[On the count]
The taxonomy contains \textbf{27} skills across the four categories, and the app
tracks and displays all 27. (How each one is scored is covered in
\S\ref{sec:scoring}.)
\end{notebox}

\subsection{The 7 player levels}

Your level is derived only from the rating carried in your games' PGN headers
(Chess.com and Lichess exports include it). Upload unrated games and the app
honestly shows no level rather than guessing one.

\begin{center}
\tabnum
\begin{tabularx}{\linewidth}{@{}c l l Y@{}}
\toprule
\textbf{ID} & \textbf{Name} & \textbf{Rating} & \textbf{Dominant failure mode} \\
\midrule
L1 & Newcomer & 0--799 & Tactical invisibility --- misses what was just threatened \\
L2 & Beginner & 800--1199 & Pattern blindness --- misses tactics beyond one move \\
L3 & Casual club player & 1200--1499 & Contextual blindness --- misses in game what they solve in puzzles \\
L4 & Improving club player & 1500--1799 & Strategic vacuum --- sound moves, no unifying plan \\
L5 & Strong club player & 1800--1999 & Reactive play --- doesn't prevent the opponent's ideas \\
L6 & Expert / CM & 2000--2299 & Conversion failure --- wins strategically, can't finish \\
L7 & Master / Titled & 2300+ & Highly individual --- needs personalised diagnosis \\
\bottomrule
\end{tabularx}
\end{center}

\subsection{The 5 plateaus}

A plateau is a named stall pattern, auto-diagnosed from the \emph{shape} of your
skill vector --- never self-reported. Zones deliberately overlap; ratings above
2200 have no formulaic plateau by design (that band is individual-coaching
territory).

\begin{center}
\tabnum
\begin{tabularx}{\linewidth}{@{}l l Y@{}}
\toprule
\textbf{Plateau} & \textbf{Rating zone} & \textbf{Target skills} \\
\midrule
The blunder wall & 0--1399 & tactical pattern recognition, tactical consistency \\
The strategy desert & 1400--1600 & strategic planning, pawn endings, endgame principles \\
The conversion ceiling & 1500--2000 & converting advantages, rook endings, endgame precision \\
The prophylaxis gap & 1700--1900 & prophylaxis, defence \& counterplay \\
The precision boundary & 2000--2200 & calculation precision, endgame precision, time mgmt, resilience \\
\bottomrule
\end{tabularx}
\end{center}

\begin{rulebox}{When a plateau is diagnosed}
A plateau is only named when \textbf{(a)} your rating falls in its zone,
\textbf{(b)} at least one of its target skills has real evidence, and
\textbf{(c)} those target skills together carry at least \textbf{5} evidence
samples. Among all qualifying plateaus, the one with the \emph{lowest average
mastery} across its evidenced targets wins. A brand-new account is given no
diagnosis at all.
\end{rulebox}

% ===========================================================================
\section{The Player's Journey}
% ===========================================================================

This section is the ``how to play'' walkthrough: the screens, in the order you
meet them. The left sidebar groups them so the navigation itself tells you where
you are: \textbf{Play} (Play vs AI) $\rightarrow$ \textbf{Your games} (Games,
Library) $\rightarrow$ \textbf{Diagnosis} (Progress, Player Model) $\rightarrow$
\textbf{Training} (Training Plan, Drills).

\begin{notebox}[First time in? Take the tour]
On your first sign-in a short \textbf{guided tour} walks you through the whole
loop, and a \textbf{``Try a demo game''} shortcut analyses a famous master game
so you see a real review in seconds --- before you've uploaded anything. Both
live behind the floating \textbf{?} button (bottom-right), which also opens the
AI help assistant. See \S\ref{sec:help}.
\end{notebox}

\subsection{Register and sign in}

Registration asks for a display name, an email, and a password of at least 8
characters. New accounts are created as role \emph{player} on the \emph{free}
tier. Sign-in is rate-limited (20 attempts per 15 minutes per address) and
returns a session token stored in your browser; it slides forward for 14 days
of activity.

\subsection{Play vs AI}

\emph{No game to upload yet? Play one.} Pick your colour and a difficulty from
\textbf{Beginner} to \textbf{Max}, then play on an interactive board --- click a
piece, then a highlighted square. Your opponent is \textbf{Stockfish}, never the
language model: at lower levels it searches shallower and sometimes picks a
weaker (but always legal) move; at Max it plays the engine's best every time.
A \textbf{Coach} panel offers up to three graded hints for the position you're
facing --- the same engine-grounded, can't-hallucinate hints as Drills, so you
get real in-match help without the model ever choosing a move for you.

When the game ends, \textbf{Save \& analyse} runs it through the same pipeline as
an uploaded game: it becomes one of \emph{your} games, so it feeds your player
model and generates drills (the opponent is labelled ``Stockfish (Level~N)'' so
practice games are easy to spot). This is the way to bootstrap a profile with no
games to import --- and, unlike a loaded library classic, a game you played
\emph{is} yours, so it counts.

\subsection{The Dashboard (``Games'')}

Your landing screen. Three stat cards summarise everything analysed so far ---
\textbf{Games analysed}, \textbf{Average accuracy}, and \textbf{Games with a
blunder} --- above a list of recent games. Each row shows the opening (ECO),
both players with your side highlighted, the colour-coded result, time control,
length, a blunder tag if any, and your accuracy. Accuracy is coloured green at
85+, gold at 75+, red below.

\subsection{Upload games}

Choose the source (Chess.com export, Lichess export, or manual PGN), tell the
app \textbf{your username as it appears in the PGN} so it can detect which
colour you played, then paste PGN text or drop a \code{.pgn} file. On submit,
each game streams its analysis progress live. Unparseable games are listed with
a reason and never sink the batch.

\begin{rulebox}{Upload limits \& quotas}
At most \textbf{50 games} per upload. Re-uploads of a game you already have are
detected by content hash and cost nothing. Analysis is metered by tier:
anonymous \textbf{1/day}, free \textbf{5/day}, Pro and Academy \textbf{unlimited}.
Duplicate and over-quota games are never charged.
\end{rulebox}

\subsection{Game Review}

The move-by-move board. On the left, an evaluation bar, the board oriented to
your colour, and navigation (buttons or keyboard arrows / Home / End). Every
move you step to shows a \textbf{classification badge}, the centipawns lost, how
long you thought, the engine's preferred move, and the top three engine lines.
On the right: live analysis progress, a \textbf{Game accuracy} card (White vs
Black, with the engine depth and a tally of each classification), the full move
list, and an \textbf{Engine summary} coaching note. \emph{This note is the first
place you meet the AI:} a language model writes it in plain English --- but only
from verified facts (your result, accuracy, and worst move), and it is re-checked
by the claim guard before you see it, so it is badged ``engine-grounded''. If the
model is unavailable or its draft fails the check, you get a deterministic,
engine-only summary badged ``deterministic'' instead. Either way, no sentence can
contradict the engine (\S\ref{sec:ai}).

\subsection{Player Model}

Headed \emph{``Every score has a receipt. No black box.''} On the left: your
current level, your diagnosed plateau, an eight-axis skill radar, and a
mastery-over-time sparkline. On the right: the full skill list with mastery
bars and trend arrows, and an \textbf{Evidence} panel. Click any skill with
evidence and it loads up to 20 receipts --- the real moves behind the number,
each marked as supporting the skill (\yes) or counting against it
(\no).

\subsection{Progress --- your guide to success}

Headed \emph{``Where you are, what you've achieved, and the single best thing to
do next.''} This is the screen to open whenever you're unsure what to work on. It
shows three things:
\begin{itemize}[leftmargin=1.4em, itemsep=1pt]
  \item a \textbf{Coach's summary} --- \emph{AI-written}, in plain language,
  grounded strictly in your real numbers (games analysed, evidenced skills,
  level, plateau, drill retention). It never invents a statistic; if the model
  is down you still get a deterministic summary of the same facts.
  \item \textbf{Your next move} --- a single, concrete next action chosen by a
  fixed rule, not the AI: upload your first game $\rightarrow$ analyse enough to
  unlock a plan $\rightarrow$ clear today's drills $\rightarrow$ study your
  lowest well-evidenced skill. A \textbf{Go} button takes you straight there.
  \item \textbf{Key stats} and a \textbf{What you've done} list --- deterministic,
  fact-by-fact, so the encouragement is always earned.
\end{itemize}

\subsection{Training Plan}

Headed \emph{``Prescribed from your diagnosis, not a generic curriculum.''} It
opens with \textbf{The hypothesis} --- a sentence tied to your diagnosed plateau
--- then up to three \textbf{focus blocks}. Each names a weak skill, its mastery
and evidence count, a rationale, and \textbf{matched reading} drawn from a book
corpus mapped to the taxonomy. A skill needs at least \textbf{3} evidence points
before it is allowed into a plan.

\subsection{Drills}

Headed \emph{``Every position here is one you actually reached and misplayed.''}
Each drill is one of your own mistake positions presented as a four-way
multiple choice --- the engine's best move plus three legal decoys. Answer
correctly and the card is scheduled further out; answer wrongly and it returns
tomorrow. A stat block tracks cards \textbf{due}, your day \textbf{streak}, and
\textbf{retention}. The scheduler is SM-2 spaced repetition (\S\ref{sec:mechanics}).

\textbf{Stuck? Ask the Coach.} A \textbf{Coach} panel offers up to three
\emph{graded, Socratic} hints: first a nudge to look at forcing moves, then the
\emph{kind} of move (a capture, a check, a quiet improvement), and only last the
move itself. Crucially, these hints are \emph{engine-grounded, not AI guesswork}
--- each is derived directly from the engine's own best move, so a hint can
never be wrong. Taking a hint is honest about what it means: a hinted-correct
counts as \textbf{partial credit}, so the card comes back sooner than one you
solved unaided (\S\ref{sec:mechanics}). Hints and help chats are recorded so
your effort is transparent.

\subsection{Library}

Two tabs. \textbf{Games} is a searchable catalogue of master games --- the
golden-start library ships with roughly \textbf{6{,}000} games from Morphy,
Capablanca, Fischer, Tal, and Kasparov, and grows via \code{npm run
library:import}. You can \textbf{search} by player, opening, or event,
\textbf{filter} by ECO code, result, and source, \textbf{sort}, and page through
the results. Loading a game runs full engine analysis on it for review (one
quota unit); because a historical game has no honest ``your side,'' a loaded
classic is \emph{not} folded into your player model --- studying Morphy teaches
you about Morphy, not about you. \textbf{Opening explorer} is an interactive
tree: at every position it shows what the master library played there against
what \emph{you} have played there in your own analysed games --- a direct read
on where your repertoire leaks.

\subsection{Account \& data}

GDPR controls. \textbf{Export} downloads every row keyed to you --- games,
moves, analyses, skill scores, evidence, coaching links --- as a single JSON
bundle (password fields stripped). \textbf{Delete} requires typing
\code{DELETE}, then transactionally erases your data, tombstones the account,
and logs you out. The deletion itself is recorded in an audit log for
compliance.

\subsection{Help, anytime}\label{sec:help}

The floating \textbf{?} button (bottom-right, on every screen) opens three ways
to get unstuck:
\begin{itemize}[leftmargin=1.4em, itemsep=1pt]
  \item \textbf{Take the guided tour} --- replays the eight-step walkthrough of
  the whole loop; it auto-runs once on your first visit and is always skippable.
  \item \textbf{Try a demo game} --- loads a famous master game and analyses it
  immediately, with a short note explaining what you're watching, so you can see
  a full review before uploading anything of your own.
  \item \textbf{Ask the help assistant} --- a chat window backed by AI. Ask
  ``how do I\dots'' or ``what does this mean'' in plain language and it explains,
  referring you to the right screen. It answers \emph{only} from a curated
  knowledge base about how Master Chess works, plus your own account context; it
  will not invent features or chess advice, and it hands you off to the Account
  page (or your administrator) for anything out of scope (\S\ref{sec:ai}).
\end{itemize}

% ===========================================================================
\section{The Mechanics}\label{sec:mechanics}
% ===========================================================================

This is the rules reference: exactly how a game becomes a diagnosis. The
worked example in \S\ref{sec:worked} then runs all of it on one real game.

\subsection{The pipeline, end to end}

\begin{enumerate}[leftmargin=1.4em, itemsep=2pt]
  \item \textbf{Ingest.} The PGN is split (comment-safe), loaded move by move,
  and stored with before/after positions, clock times, and detected player
  colour.
  \item \textbf{Evaluate.} Every position is scored by Stockfish (results are
  cached, since openings recur), producing an evaluation and ranked candidate
  lines for each.
  \item \textbf{Classify.} Each move is labelled \best\ \good\ \inacc\ \mistake\ \blun\
  by how much it lost against the engine's best, scaled by game phase.
  \item \textbf{Infer evidence.} Fixed rules read the classified moves and emit
  ``for'' / ``against'' evidence for specific skills.
  \item \textbf{Update the model.} Evidence is folded into each skill's
  Bayesian Knowledge Tracing state; a snapshot records level and plateau.
  \item \textbf{Harvest \& prescribe.} Mistakes become drills; the training
  plan is rebuilt; finally the narrator writes a claim-guarded summary.
\end{enumerate}

\begin{notebox}[Failure isolation]
The engine analysis is committed and shown to you \emph{before} the enrichment
step (evidence, model, drills, narration) runs. Enrichment is wrapped so that
if any part of it fails, your analysis still stands as ``done'' --- and it is
idempotent, so a restart mid-way never double-counts evidence.
\end{notebox}

\subsection{Move classification}

A move's raw loss is how many centipawns worse it is than the engine's best
reply. Thresholds are \textbf{phase-scaled} --- the endgame is judged more
strictly than the opening, where many moves are near-equal:

\begin{center}
\tabnum
\begin{tabularx}{\linewidth}{@{}l c c c c@{}}
\toprule
\textbf{Phase} & \textbf{good $>$} & \textbf{inaccuracy $\geq$} & \textbf{mistake $\geq$} & \textbf{blunder $\geq$} \\
\midrule
Opening & 5 & 60 & 150 & 350 \\
Middlegame & 5 & 50 & 120 & 300 \\
Endgame & 5 & 40 & 90 & 200 \\
\bottomrule
\end{tabularx}\\[4pt]
{\footnotesize\color{slate}Values are centipawns of \emph{effective} loss. At or below ``good'', a move is \best.}
\end{center}

\subsubsection{Phase detection is queen-aware}
Phase is not just move number. Very low material (a lone minor per side or less,
$\leq 500$cp) forces \textbf{endgame} even early --- an early queen trade into a
king-and-pawn position is an endgame. Otherwise the first 20 plies are
\textbf{opening}; after that, low non-pawn material \emph{with no queens on the
board} ($\leq 1400$cp) is \textbf{endgame}; everything else is
\textbf{middlegame}. Queens keep a position out of the endgame even when other
material is thin.

\subsubsection{Win-probability dampening}
Raw centipawns mislead in decided positions: going from +900 to +700 is not a
blunder, because you were completely winning before and after. The classifier
converts evaluations to win probability on the Lichess curve and, \emph{only
when both endpoints sit deep in a decided zone}, blends the raw loss toward the
much smaller win-probability loss. A genuine thrown win --- say +2000 collapsing
to +200 --- is \emph{not} dampened, because the second endpoint is no longer
decided. Missing a forced mate is always at least a mistake regardless of
centipawns.

\begin{rulebox}{Accuracy score}
Per move, the app computes win-probability points lost, maps each to a 0--100
accuracy on the Lichess curve
($103.1668\,e^{-0.04354x}-3.1669$), then blends the arithmetic and harmonic
means of the whole game. The harmonic term is what makes a single game-losing
blunder cost more than a plain average would.
\end{rulebox}

\subsection{Turning moves into evidence}

Rules fire only on \emph{your} moves and emit evidence with a direction and a
weight in $[0,1]$. Error severity is graded blunder $=1$, mistake $=0.6$,
inaccuracy $=0.3$. A representative selection:

\begin{center}
\tabnum
\begin{tabularx}{\linewidth}{@{}l Y c c@{}}
\toprule
\textbf{Rule} & \textbf{Fires when} & \textbf{Skill} & \textbf{Dir / weight} \\
\midrule
missed-forced-mate & you had a forced mate and gave it up & calculation precision & against / 1.0 \\
opening-error & opening-phase mistake or blunder & opening principles & against / sev. \\
middlegame-blunder & a blunder in the middlegame & tactical consistency & against / 1.0 \\
\dots{}-blunder-pattern & (same move, also) & tactical pattern recog. & against / 0.5 \\
hanging-piece-left & error that leaves a piece capturable \& undefended & piece activity & against / 0.4 \\
endgame-error & endgame mistake or blunder & endgame precision & against / sev. \\
snap-blunder & error played in $\leq 3$s & time management & against / 0.4 \\
conversion-failure & reached $\geq +400$cp but didn't win & converting advantages & against / 0.8 \\
sustained-initiative & grew an even game to $\geq +500$ cleanly & attack \& initiative & for / 0.55 \\
held-worse-position & were $\leq -300$ yet didn't lose & defence \& counterplay & for / 0.5 \\
consistent-repertoire & same opening, played cleanly, $\geq 3$ games & opening repertoire & for / 0.4 \\
sound-novelty & sound engine-approved deviation from book & opening prep. & for / 0.3 \\
prophylaxis probe & a null-move probe reveals a real threat & prophylaxis & for/against \\
\bottomrule
\end{tabularx}
\end{center}

\begin{notebox}[The prophylaxis probe]
To judge prophylaxis, the app performs a ``null-move'' probe: it lets your
opponent move twice and asks the engine how much that would have cost you. If a
real threat existed ($\geq 300$cp) and you let it land, that is evidence against
your prophylaxis; if it existed and you neutralised it, that is evidence for.
The weight is derived from the threat's size, never from the language model.
Probes are skipped when you were in check (the position would be illegal) and
capped at six per game.
\end{notebox}

\subsection{How each skill is scored}\label{sec:scoring}

Every skill's mastery comes from the same chain --- \emph{engine classifies your
moves} $\rightarrow$ \emph{a fixed rule turns specific patterns into for/against
evidence} $\rightarrow$ \emph{Bayesian Knowledge Tracing folds that evidence into
a 0--100 number} (next subsection). What differs between skills is the
\textbf{rule}: what, exactly, is measured. The Player Model shows this method in
plain language for each skill, so a score is never a mystery. Here is the full
map.

\begin{center}
\tabnum
\begin{xltabular}{\linewidth}{@{}l Y@{}}
\toprule
\textbf{Skill} & \textbf{What is measured} \\
\midrule
\endhead
\multicolumn{2}{@{}l}{\textbf{\color{forest}OPENING}}\\
Opening principles & Errors in the first 20 plies count against; clean opening play counts lightly for. \\
Opening repertoire & Replaying the same ECO line in 3+ games with low opening cp-loss counts for. Variety is never penalised. \\
Opening prep \& novelties & Leaving the master-game book (moves 4--10) with an engine-approved move counts for (sound novelties only). \\
\addlinespace
\multicolumn{2}{@{}l}{\textbf{\color{forest}MIDDLEGAME}}\\
Tactical consistency & A middlegame blunder counts against --- the ``did you blunder-check?'' measure. \\
Tactical pattern recognition & Middlegame mistakes/blunders where the engine clearly saw better count against. \\
Calculation precision & Missing a forced mate the engine had found counts heavily against. \\
Piece activity \& coordination & Leaving your own piece attacked and undefended (on an error) counts against --- a real board check, not just cp-loss. \\
Attack \& initiative & Growing an even position to clearly winning ($+5$) counts for. \\
Defence \& counterplay & Being clearly worse ($-3$) yet drawing or winning counts for. A loss is not counted against. \\
Converting advantages & Reaching a won position ($+400$cp) but not winning counts strongly against. \\
Prophylaxis & A null-move engine probe gauges a position's standing threat; a real threat you let land counts against, one you neutralised counts for. \\
\addlinespace
\multicolumn{2}{@{}l}{\textbf{\color{forest}ENDGAME}}\\
Endgame principles & Endgame mistakes/blunders count against; sound endgame play counts lightly for. \\
Endgame precision \& conversion & Endgame cp-loss counts against --- precision once the board simplifies. \\
Pawn / rook / knight / bishop endings & The same endgame errors are \emph{also} attributed to the specific piece-type ending they occurred in. \\
\addlinespace
\multicolumn{2}{@{}l}{\textbf{\color{forest}PSYCHOLOGY \& MENTAL}}\\
Time management & A blunder/mistake played in under 3 seconds counts against. \\
Thought process \& candidate moves & A blunder/mistake after 90+ seconds of thought counts against (had time, still missed it). \\
Psychological resilience & The spread of your cp-losses and whether errors spiked after your first blunder (``tilt'') vs. stayed controlled. \\
\bottomrule
\end{xltabular}
\end{center}

\begin{rulebox}{``Not yet scored'' is a real answer}
Seven of the 27 skills have no honest rule yet, because they need judgment
beyond move classification: \emph{strategic planning}, \emph{positional play},
\emph{patience \& long-term thinking}, \emph{game transition}, \emph{pattern-
recognition speed}, \emph{intuition \& practical judgment}, and \emph{psychological
warfare}. Rather than show a made-up number, the app marks these \textbf{Not yet
scored}. So a blank skill means ``not measured yet'' --- never ``you are bad at
this.''
\end{rulebox}

\subsection{Bayesian Knowledge Tracing: how a score moves}

Each skill holds a probability \code{pKnow} that you have mastered it; mastery
shown on screen is simply $\text{round}(\mathtt{pKnow}\times 100)$. A new skill
starts at \code{pKnow}~$=0.3$ (mastery 30). Each piece of evidence updates it
with a weighted Bayesian step, using fixed parameters (slip $0.1$, guess $0.2$,
learning transition $0.15$).

\begin{rulebox}{The update, in words}
A ``for'' observation raises \code{pKnow} toward the Bayesian posterior and adds
a small learning bump; an ``against'' observation lowers it. \textbf{Weight
scales the whole step}: a weight-1 blunder moves the number hard, a weight-0.3
inaccuracy nudges it, a weight-0 observation does nothing at all. Evidence is
folded in sequence, so order within a skill matters.
\end{rulebox}

The concrete effect of a \emph{single} full-weight observation from a fresh
skill (mastery 30):

\begin{center}
\tabnum
\begin{tabularx}{\linewidth}{@{}Y c c@{}}
\toprule
\textbf{Observation} & \textbf{New \code{pKnow}} & \textbf{New mastery} \\
\midrule
one ``for'', weight 1.0 & 0.710 & \textbf{71} \\
one ``against'', weight 1.0 & 0.051 & \textbf{5} \\
one ``against'', weight 0.3 & 0.225 & \textbf{23} \\
weight 0 (any) & 0.300 & 30 (no change) \\
\bottomrule
\end{tabularx}
\end{center}

Confidence in the whole model grows with the number of samples, capped:
$\min(95,\ \text{samples}\times 3)$ percent.

\subsection{Own-mistake drills and SM-2 scheduling}

Every mistake or blunder of yours that lost at least \textbf{100}cp --- and for
which the engine has a genuinely better move --- becomes a drill. The correct
answer is the engine's own best move; positions are de-duplicated so you never
drill the same one twice.

Scheduling is SM-2-lite. A card carries an \emph{ease} factor (start 2.0,
floored 1.3, capped 2.6) and an interval that grows multiplicatively so a
correct answer can never shorten it.

\begin{rulebox}{The schedule}
\textbf{Correct:} ease rises (more for better-known skills), and the next
interval is $\text{interval}\times\text{ease}$, capped at \textbf{21 days}.
\textbf{Incorrect:} ease drops by 0.2, the card returns in 5 minutes, and a
lapse is recorded. After \textbf{4 lapses} a card is suspended as a ``leech''
and stops being scheduled. A card only counts toward your streak and retention
when it was genuinely due --- you cannot grind the same card for credit.
\textbf{Hinted-correct} (you used the Coach) is partial credit: no ease reward,
no streak, and only a gentle interval bump, so the card returns sooner than an
unaided solve.
\end{rulebox}

An answer is also accepted if it is any engine move within about 25cp of the
best --- two equally winning moves are both marked correct.

\begin{center}
\tabnum
\begin{tabularx}{\linewidth}{@{}c c c Y@{}}
\toprule
\textbf{Correct answer \#} & \textbf{Ease} & \textbf{Next interval} & \textbf{Meaning} \\
\midrule
1 & 2.08 & 2.08 days & graduates from ``new'' \\
2 & 2.16 & $\approx$4.5 days & \\
3 & 2.24 & $\approx$10 days & \\
4 & 2.32 & 21 days (capped) & considered learned \\
\bottomrule
\end{tabularx}\\[3pt]
{\footnotesize\color{slate}Progression for a card on a mastery-30 skill, answered correctly each time.}
\end{center}

% ===========================================================================
\section{A Complete Game of Play}\label{sec:worked}
% ===========================================================================

We now run the whole system on one real game: Paul Morphy's ``Opera Game''
(Paris, 1858), analysed as though \emph{you uploaded it and played Black} --- the
allied Duke and Count. Every number, verdict, and skill change in the boxes
below is the \textbf{actual output} of the Master Chess pipeline run on this
game's moves, not an illustration.

\begin{notebox}[This example was generated, not hand-written]
The figures here come from running the app's real classifier, evidence rules,
and Bayesian update on the stored analysis of this game. One consequence worth
stating up front: the model's job is to read the moves honestly, and honest
reading sometimes surprises --- as it does here.
\end{notebox}

\begin{playbox}[Step 1 — Upload \& ingest]
You upload the game as Black. The ingester reads the header ---
\code{White: Paul Morphy}, \code{Black: Duke Karl / Count Isouard},
\code{Result: 1-0}, \code{ECO: C41}, \code{Opening: Philidor Defense} --- matches
your name to the Black seat, and replays all 33 plies, storing the position
before and after each. The moves:

\vspace{4pt}
{\ttfamily\small 1.e4 e5 2.Nf3 d6 3.d4 Bg4 4.dxe5 Bxf3 5.Qxf3 dxe5 6.Bc4 Nf6
7.Qb3 Qe7 8.Nc3 c6 9.Bg5 b5 10.Nxb5 cxb5 11.Bxb5+ Nbd7 12.O-O-O Rd8 13.Rxd7
Rxd7 14.Rd1 Qe6 15.Bxd7+ Nxd7 16.Qb8+ Nxb8 17.Rd8\#}

\vspace{4pt}
\footnotesize\color{slate}
A historical game carries no rating header, so --- as \S\ref{sec:mechanics}
requires --- no \emph{level} or \emph{plateau} is diagnosed from it. The skill
evidence below is generated all the same.
\end{playbox}

\begin{playbox}[Step 2 — Evaluate \& classify (all of Black's moves)]
Stockfish scores every position; the classifier judges each Black move against
the best available, phase-scaled. This is the complete, verbatim classification:

\begin{center}
\tabnum
\begin{tabularx}{\linewidth}{@{}l l l c Y@{}}
\toprule
\textbf{Move} & \textbf{Played} & \textbf{Phase} & \textbf{cp lost} & \textbf{Verdict} \\
\midrule
1\dots e5 & & opening & 3 & \best \\
2\dots d6 & & opening & 9 & \good \\
3\dots Bg4 & & opening & 65 & \inacc \\
4\dots Bxf3 & & opening & 54 & \good \\
5\dots dxe5 & & opening & 0 & \best \\
6\dots Nf6 & & opening & 67 & \inacc \\
7\dots Qe7 & & opening & 18 & \good \\
8\dots c6 & & opening & 23 & \good \\
9\dots b5 & & opening & 84 & \inacc \\
10\dots cxb5 & & opening & 139 & \inacc \\
11\dots Nbd7 & & middlegame & 33 & \good \\
12\dots Rd8 & & middlegame & 58 & \inacc \\
13\dots Rxd7 & & middlegame & 1 & \best \\
14\dots Qe6 & & middlegame & 132 & \mistake \\
15\dots Nxd7 & & middlegame & mate & \blun \\
16\dots Nxb8 & & middlegame & 0 & \best \\
\bottomrule
\end{tabularx}
\end{center}

Two things the real data teaches. First, the game is classified
\textbf{opening} then \textbf{middlegame} throughout --- the queen-aware phase
rule keeps it out of the endgame despite the piece sacrifices. Second, the
famous ``losing'' move \textbf{9\dots b5 is only an inaccuracy} (84cp). The move
the engine actually flags as the \textbf{blunder is 15\dots Nxd7}, which walked
into \code{16.Qb8+ Nxb8 17.Rd8\#} --- a forced mate (shown as ``mate'' rather
than a centipawn figure).
\end{playbox}

\begin{playbox}[Step 3 — Inferred evidence (exactly what the rules emit)]
The rule engine reads those classifications. A crucial detail visible here:
\textbf{inaccuracies emit no per-move skill evidence} --- only clearly good
moves (\best/\good) and clear errors (\mistake/\blun) do. So Bg4, Nf6, b5, cxb5,
and Rd8 produce nothing. What actually fires:

\begin{center}
\tabnum
\begin{tabularx}{\linewidth}{@{}l Y l@{}}
\toprule
\textbf{Skill} & \textbf{From (real moves)} & \textbf{Dir / wt} \\
\midrule
Opening principles & e5, d6, Bxf3, dxe5, Qe7, c6 (six sound opening moves) & for / 0.15 \\
Tactical consistency & Nbd7, Rxd7, Nxb8 (sound middlegame) & for / 0.12 \\
Tactical consistency & 15\dots Nxd7 (middlegame blunder) & against / 1.0 \\
Tactical pattern recog. & 14\dots Qe6 (middlegame mistake) & against / 0.6 \\
Tactical pattern recog. & 15\dots Nxd7 (blunder pattern) & against / 0.5 \\
\bottomrule
\end{tabularx}
\end{center}

Note what does \emph{not} fire: no ``conversion failure'' (Black never reached
$+400$), no ``held-worse-position'' credit (Black lost), and no resilience
verdict (the run did not meet the sample conditions). The model records only
what the game shows.
\end{playbox}

\begin{playbox}[Step 4 — The player model updates (real Bayesian output)]
Each skill folds its evidence in sequence from a starting mastery of 30. The
actual results:

\begin{center}
\tabnum
\begin{tabularx}{\linewidth}{@{}Y c c c@{}}
\toprule
\textbf{Skill} & \textbf{Net evidence} & \textbf{Mastery} & \textbf{Trend} \\
\midrule
Opening principles & 6 for & $30 \to 66$ & $\uparrow$ \\
Tactical consistency & 3 for, 1 against (wt 1) & $30 \to 12$ & $\downarrow$ \\
Tactical pattern recognition & 2 against & $30 \to 9$ & $\downarrow$ \\
\bottomrule
\end{tabularx}
\end{center}

The surprise is honest and instructive: \textbf{opening principles rises to 66}
even though Black \emph{lost the game}. By cp-loss, Black's opening was sound ---
the defeat came from two middlegame errors, not the opening. A human eye
remembers ``the game where I got mated''; the model measures each phase on its
own evidence. That is precisely the black-box the receipts exist to open.
\end{playbox}

\begin{playbox}[Step 5 — Drills harvested]
Two moves qualify as own-mistake drills --- classified \mistake\ or \blun\ and
losing at least 100cp:

\begin{center}
\tabnum
\begin{tabularx}{\linewidth}{@{}l l Y@{}}
\toprule
\textbf{Position before} & \textbf{Skill} & \textbf{Correct answer} \\
\midrule
14\dots Qe6 (132cp) & tactical pattern recognition & the engine's best move \\
15\dots Nxd7 (mate) & tactical consistency & the engine's best move \\
\bottomrule
\end{tabularx}
\end{center}

Each becomes a four-choice card, due immediately, entering SM-2 scheduling.
Note that \code{9\dots b5} and \code{10\dots cxb5} are \emph{not} harvested ---
b5 lost under 100cp, and both were inaccuracies rather than outright errors.
\end{playbox}

\begin{playbox}[Step 6 — Training plan prescribed]
A skill needs at least 3 evidence samples to enter a plan. Here that admits
\emph{opening principles} (6 samples) and \emph{tactical consistency} (4); with
only 2 samples, \emph{tactical pattern recognition} is honestly held back
despite being the lowest score. Because the historical game has no rating, no
plateau is diagnosed, so the plan uses its no-plateau hypothesis:

\vspace{3pt}
\textbf{The hypothesis.} \emph{``Your weakest evidenced skills right now are
tactical consistency, opening principles. No plateau pattern is diagnosable yet
(needs a PGN-header rating), but these are the lowest-mastery skills we have real
evidence for.''}

\vspace{4pt}
\textbf{Focus block 1 --- Tactical consistency} (mastery 12, 4 samples).
\emph{Matched reading} from the corpus: \emph{Art of Combination} (Blokh),
\emph{Encyclopedia of Chess Combinations} (Chess Informant), \emph{1000
Checkmate Combinations} (Henkin).

\vspace{2pt}
\textbf{Focus block 2 --- Opening principles} (mastery 66, 6 samples), carrying
its own matched titles. It ranks second only because it is the next evidenced
skill --- the plan is transparent about that rather than inventing urgency.
\end{playbox}

\begin{playbox}[Step 7 — The engine-grounded note]
Finally the narrator is handed \emph{only} verified facts --- result, accuracy,
blunder count, and the worst move (\code{15\dots Nxd7}, a mate-allowing blunder)
--- and asked for two to four sentences, then re-checked by the claim guard. A
passing note reads like:

\vspace{2pt}
\emph{``Your opening was solid --- it was the middlegame that cost you. 14\dots
Qe6 let the pressure build, and 15\dots Nxd7 walked into a forced mate with
16.Qb8+ and 17.Rd8\#. The habit to drill is the blunder-check before a
recapture when your king is still in the centre.''}

\vspace{4pt}
Every move and claim there maps to the real record. Had the draft said ``your
9\dots b5 was the losing blunder'' --- contradicting the classifier, which
graded b5 an inaccuracy --- the guard would reject it and you would see the
deterministic summary instead. \emph{(That is the exact error this manual's first
draft made before it was checked against the engine.)}
\end{playbox}

\begin{notebox}[What you would do next]
Open \emph{Drills} the next day: the \code{15\dots Nxd7} and \code{14\dots Qe6}
positions are waiting. Choose the engine's move and the card schedules out; miss
it and it returns tomorrow until the gap closes. That is the loop --- diagnose,
explain, prescribe --- closing on your own moves.
\end{notebox}

% ===========================================================================
\section{Reference Card}
% ===========================================================================

\subsection{Screens at a glance}

\begin{center}
\begin{tabularx}{\linewidth}{@{}l Y@{}}
\toprule
\textbf{Play vs AI} & play Stockfish at 5 levels; save \& analyse to feed your model \\
\textbf{Games} & dashboard: stats + recent games \\
\textbf{Upload} & paste/drop PGN, live analysis \\
\textbf{Game Review} & per-move engine analysis, accuracy, AI summary \\
\textbf{Progress} & AI coach's summary + your single next move \\
\textbf{Player Model} & level, plateau, 27 skills, evidence receipts \\
\textbf{Training Plan} & hypothesis + 3 focus blocks + books \\
\textbf{Drills} & SM-2 review of your own mistakes \\
\textbf{Library} & searchable master games + opening explorer \\
\textbf{Account} & GDPR export / delete \\
\textbf{? (help)} & guided tour, demo game, AI help chat \\
\bottomrule
\end{tabularx}
\end{center}

\subsection{Numbers worth remembering}

\begin{center}
\tabnum
\begin{tabularx}{\linewidth}{@{}Y l@{}}
\toprule
Skills / categories / levels / plateaus & 27 / 4 / 7 / 5 \\
Fresh skill mastery & 30 \\
Blunder threshold (mid / open / end) & 300 / 350 / 200 cp \\
Drill harvest threshold & 100 cp lost \\
Drill accept tolerance & $\sim$25 cp of best \\
SM-2 ease (start / floor / cap) & 2.0 / 1.3 / 2.6 \\
Max review interval & 21 days \\
Leech suspension & 4 lapses \\
Plateau minimum samples & 5 \\
Plan minimum samples per skill & 3 \\
Free-tier analyses / day & 5 \\
Max games per upload & 50 \\
\bottomrule
\end{tabularx}
\end{center}

\subsection{Classification legend}

\begin{center}
\best\quad\good\quad\inacc\quad\mistake\quad\blun
\end{center}

\vfill
\begin{center}
{\color{rulehair}\rule{0.5\linewidth}{0.6pt}}\\[4pt]
{\sffamily\footnotesize\color{slate}
Master Chess \textbullet\ Field Manual \textbullet\ Diagnose \textbullet\ Explain
\textbullet\ Prescribe}
\end{center}

\end{document}
